\documentclass[11pt]{article}
\usepackage{fontspec,amsmath,amssymb,geometry}
\setmainfont{Times New Roman}\geometry{margin=0.7in}
\newcommand{\topic}[1]{\par\leftskip=0pt\section*{\normalfont\large #1}}\newcommand{\initial}[1]{{\Large #1}}\newcommand{\result}[1]{\par\leftskip=0pt\medskip\noindent\hspace*{1em}{\bfseries #1}\par\nobreak\leftskip=1em\noindent}\renewcommand{\textbf}[1]{\result{#1}}
\newcommand{\question}[1]{\par\leftskip=0pt\medskip\noindent{\bfseries Q#1.}\par\nobreak\leftskip=1em\noindent}
\newcommand{\subquestion}[2]{\par\noindent\hangindent=2em\hangafter=1\makebox[1.5em][l]{(#1)}\hspace{0.5em}#2\par}
\newenvironment{chaptermap}{\begin{center}\begin{minipage}{0.82\textwidth}\small}{\end{minipage}\end{center}\medskip}
\title{Determinants}\author{T.T.}\date{}
\begin{document}\maketitle
\begin{chaptermap}\itshape Determinants assign a signed volume and test invertibility. Their defining properties yield row-operation rules, cofactor expansions, permutation formulas, and inverse formulas.\end{chaptermap}
\topic{\initial{D}EFINITION AND \initial{P}ROPERTIES}
\[\det I=1,\qquad \det(AB)=\det(A)\det(B),\qquad A\text{ invertible}\Longleftrightarrow\det A\ne0.\]
\textbf{Row rules:} swapping rows changes the sign; multiplying a row by $c$ multiplies the determinant by $c$; adding a multiple of one row to another leaves it unchanged.
\textbf{Linearity:} the determinant is linear in each row (or column) separately and alternating in its rows.
\[\det({A}^{T})=\det A,\qquad \det({A}^{-1})={\det A}^{-1},\qquad \det(cA)=c^n\det A.\]
\topic{\initial{C}OMPUTATIONAL \initial{F}ORMULAE}
\textbf{Triangular matrices:} $\det A=\prod_{i=1}^n a_{ii}$.
\[\det A=\sum_{j=1}^n a_{ij}C_{ij}=\sum_{i=1}^n a_{ij}C_{ij},\quad C_{ij}={(-1)}^{i+j}\det(A_{ij}).\]
\textbf{Cofactor formula:} $\displaystyle {A}^{-1}=\frac{1}{\det A}{\operatorname{Cof}(A)}^{T}$.
\[\det A=\sum_{\sigma\in S_n}\operatorname{sgn}(\sigma)\prod_{i=1}^n a_{i,\sigma(i)}.\]
\textbf{Cramer's rule:} if $Ax=b$ and $\det A\ne0$, $\displaystyle x_i=\frac{\det A_i(b)}{\det A}$.
\textbf{Vandermonde determinant:} $\displaystyle \det\!\left({[x_j^{i-1}]}_{i,j=1}^n\right)=\prod_{1\leq i<j\leq n}(x_j-x_i)$.
\topic{\initial{G}EOMETRY AND \initial{V}OLUME}
\[\operatorname{volume}\bigl(A({[0,1]}^n)\bigr)=|\det A|,\qquad \det A=0\Longleftrightarrow\text{columns lie in a lower-dimensional subspace}.\]
\newpage
\setlength{\parindent}{0pt}
\topic{\initial{P}ROBLEM \initial{S}ET}

\question{1} Use row operations to compute
\[
\det\begin{bmatrix}
1&2&3\\
2&5&7\\
1&1&1
\end{bmatrix}.
\]
State the determinant rule used at each step and identify the product of the pivots.

\question{2} Let $A$ and $B$ be $4\times4$ matrices with $\det A=-2$ and $\det B=3$.
\subquestion{i}{Find $\det(AB)$ and $\det({A}^{T})$.}
\subquestion{ii}{Find $\det(A^{-1})$, $\det(A^{2})$, and $\det(2A)$.}

\question{3} Let
\[
A=\begin{bmatrix}2&1&0\\1&3&2\\0&4&1\end{bmatrix}.
\]
\subquestion{i}{Find the cofactor matrix $C=(C_{ij})$ of $A$.}
\subquestion{ii}{Use a cofactor expansion to find $\det A$.}
\subquestion{iii}{Verify the identity $A{C}^{T}=(\det A)I$.}

\question{4} For
\[
A(t)=\begin{bmatrix}1&1&1\\1&t&2\\1&2&t\end{bmatrix},
\]
\subquestion{i}{Compute $\det A(t)$ by row operations.}
\subquestion{ii}{For which values of $t$ is $A(t)$ invertible?}
\subquestion{iii}{When $t=0$ or $t=2$, give a nonzero vector in the nullspace of $A(t)$.}

\question{5} Use Cramer's rule to solve
\[
2x+y=1,\qquad x+2y+z=2,\qquad y+2z=3.
\]
Compute $\det A$ and the three determinants obtained by replacing one column of $A$ by $b$.

\question{6} Use the cofactor formula to find the inverse of
\[
A=\begin{bmatrix}1&2&0\\2&1&1\\0&1&1\end{bmatrix}.
\]
Verify the result by multiplication.

\question{7}
\subquestion{i}{Find the area of the parallelogram spanned by $u=(3,2)$ and $v=(1,4)$.}
\subquestion{ii}{Find the volume of the parallelepiped spanned by
\[
u=(1,1,0),\qquad v=(1,0,1),\qquad w=(0,1,1).
\]
}

\question{8} Suppose that $A$ is an invertible matrix with integer entries.
\subquestion{i}{Prove that if $\det A=1$ or $\det A=-1$, then every entry of $A^{-1}$ is an integer.}
\subquestion{ii}{Prove the converse: if both $A$ and $A^{-1}$ have integer entries, then $\det A=1$ or $\det A=-1$.}
\end{document}
